% SOURCE_AUTHORITY: sga5_fr_workpass.tex, lines 10486--10515 (LF-normalized unit).
% SOURCE_SHA256: 791F4EFFC5E02832D5D77ED03518C8156D6F07E4C8238B03545DB93D883FBB28
Bajo las hipótesis anteriores, se ve que el conocimiento de los grupos de cohomología de $X$ con coeficientes en $\mathbb Q_\ell$ implica el conocimiento de los de $Y$, salvo en dimensión $n-1$. En estas condiciones, el cálculo de este último se reduce inmediatamente al cálculo de la característica de Euler-Poincaré de $Y$. Nos ocuparemos ahora de este cálculo en un marco algo más general.

\begin{statement}{Proposición 7.3}
Sean $X$ un $k$-esquema propio, liso, conexo, de dimensión $n$, $E$ un $\mathcal O_X$-módulo localmente libre de rango constante $r$, e $Y$ el lugar de los ceros de una sección de $V(\check E)$ transversal a la sección cero. Denotando por $p:X\to k$ el morfismo canónico y por $p_*$ el morfismo de Gysin correspondiente, se tiene la igualdad
\[
\operatorname{EP}(Y)=p_*\bigl(c(\check\Omega_X^1)c(E)^{-1}c_r(E)\bigr)
\]
en $H^0(k,\mathbb Q_\ell)=\mathbb Q_\ell$.
\end{statement}

\emph{Demostración.} Según (EGA IV 17.13.2), $Y$ es liso de dimensión pura $n-r$ y, denotando por $i:Y\to X$ la inmersión canónica, el haz conormal de $i$ es la imagen inversa por $i$ del haz conormal $\check E$ de la sección cero de $V(\check E)$. Se tiene, por tanto, una sucesión exacta
\[
0\to i^*(\check E)\to i^*(\Omega_X^1)\to\Omega_Y^1\to0,
\tag{7.3.1}
\]
de donde resulta inmediatamente la igualdad
\[
c(\check\Omega_Y^1)=i^*\bigl(c(\check\Omega_X^1)c(E)^{-1}\bigr).
\tag{7.3.2}
\]
% source check 2026-06-11: scan p322 prints reference (3.10) exactly; left unchanged.
De aquí se deduce por (3.10) la expresión siguiente para $\operatorname{EP}(Y)$:
\[
\operatorname{EP}(Y)=p_*i_*i^*\bigl(c(\check\Omega_X^1)c(E)^{-1}\bigr),
\]
es decir, teniendo en cuenta la fórmula de proyección para el morfismo $i$,
\[
\operatorname{EP}(Y)=p_*\bigl(c(\check\Omega_X^1)c(E)^{-1}i_*(1_Y)\bigr).
\]
La fórmula del enunciado se deduce inmediatamente gracias a (4.7).
